% ------------------------------------------------------------------
% preamble.tex -- hand written, book-level settings.
% This file is never touched by tools/merge_tex.py.
% Everything harvested from the source documents goes into
% preamble_generated.tex, which is loaded right after this one.
%
% Deliberately avoided: babel (no italian data files on this install),
% lmodern, biblatex, mhchem. The bibliography is a plain thebibliography.
% ------------------------------------------------------------------

\usepackage[utf8]{inputenc}
\usepackage[T1]{fontenc}
\usepackage[a4paper,inner=28mm,outer=24mm,top=25mm,bottom=28mm]{geometry}

\usepackage{amsmath,amssymb,amsfonts}
\usepackage{graphicx}
\usepackage{booktabs}
\usepackage{longtable}
\usepackage{array}
\usepackage{multirow}
\usepackage{siunitx}
\usepackage{float}
\usepackage{caption}
\usepackage{subcaption}
\usepackage{enumitem}
% expansion=false keeps microtype working on installs without scalable
% (Type 1 / cm-super) fonts -- protrusion alone is the safe subset
\usepackage[protrusion=true,expansion=false]{microtype}
\usepackage{xcolor}
\usepackage{tikz}
\usepackage{pgfplots}
\usepackage{circuitikz}
\usepackage{listings}
\usepackage{fancyhdr}
\usepackage[hidelinks,breaklinks]{hyperref}

\pgfplotsset{compat=1.18}

% figures are looked up in the chapter folders as well as ./figs
\graphicspath{{figs/}{./}{chapters/}{src/}}

% ---------------------------------------------------------------- layout
\setcounter{secnumdepth}{3}
\setcounter{tocdepth}{2}

\pagestyle{fancy}
\fancyhf{}
\fancyhead[L]{\nouppercase{\leftmark}}
\fancyhead[R]{\thepage}
\renewcommand{\headrulewidth}{0.4pt}
\fancypagestyle{plain}{\fancyhf{}\fancyfoot[C]{\thepage}%
  \renewcommand{\headrulewidth}{0pt}}

% ------------------------------------------------- chapter-opening summary
% Source documents written as articles often start with an abstract.
% merge_tex.py rewrites abstract -> chapabstract so it survives in a book.
\newenvironment{chapabstract}
  {\begin{quotation}\noindent\small\itshape}
  {\end{quotation}\par\medskip}

% ------------------------------------------------------------- code style
\lstdefinestyle{plain}{
  basicstyle=\ttfamily\footnotesize,
  keywordstyle=\bfseries,
  commentstyle=\itshape\color{gray},
  numbers=left, numberstyle=\tiny, numbersep=8pt,
  breaklines=true, frame=single, framerule=0.3pt,
  showstringspaces=false, tabsize=2, captionpos=b,
}
\lstset{style=plain}

% ------------------------------------------------------------ misc tweaks
\numberwithin{equation}{chapter}
\numberwithin{figure}{chapter}
\numberwithin{table}{chapter}
\sisetup{per-mode=symbol, range-units=single, detect-weight=true}
